%Begin


%%%%%%%%%%%%%%%%%%%%%%%%%%%%%%%%%%%%%%%%%%%%%%%%%%%%%%%%%%%%%%%
\chapter*{\S \space III. --- CLASSIFICATION DES PLONGEMENTS À CONGRUENCE PRÈS ET VOISINAGES TUBULAIRES DES 1-COMPLEXES}
\addcontentsline{toc}{section}{III. Classification des plongements à congruence près et voisinages tubulaires des 1-complexes}
\section*{}

{\bf 0.}. --- 

%%%%%%%%%%%%%%%%%%%%%%%%%%%%%%%%%%%%
\subsection*{1. Les 1-complexes circularisés.}
\addcontentsline{toc}{subsection}{1. Les 1-complexes circularisés}

{\bf 1.1}. --- 


\vskip .3cm
{\bf 1.2}. --- 


\vskip .3cm
{\bf 1.3}. --- 

\vskip .3cm
{\bf 1.4}. --- 


\vskip .3cm
{\bf 1.5}. --- 



%%%%%%%%%%%%%%%%%%%%%%%%%%%%%%%%%%%%
\subsection*{2. La catégorie des plongements comme produit 2-fibré.}
\addcontentsline{toc}{subsection}{2. La catégorie des plongements comme produit 2-fibré}

{\bf 2.1}. --- 

\vskip .3cm
{\bf 2.2}. ---


\vskip .3cm
{\bf 2.3}. ---



%%%%%%%%%%%%%%%%%%%%%%%%%%%%%%%%%%%%
\subsection*{3. Classification à congruence près des plongements de 1-complexes topologiques compacts dans les surfaces à bord généralisé orientées compactes.}
\addcontentsline{toc}{subsection}{3. Classification à congruence près des plongements de 1-complexes topologiques compacts dans les surfaces à bord généralisé orientées compactes}

{\bf 3.1}. --- 

\vskip .3cm
{\bf 3.2}. ---

\vskip .3cm
{\bf 3.3}. ---


\vskip .3cm
{\bf 3.4}. ---



%%%%%%%%%%%%%%%%%%%%%%%%%%%%%%%%%%%%
\subsection*{4. Voisinages tubulaires de 1-complexes.}
\addcontentsline{toc}{subsection}{4. Voisinages tubulaires de 1-complexes}

{\bf 4.1}. --- 

\vskip .3cm
{\bf 4.2}. --- Voisinage tubulaire associé à un plongement.

\vskip .3cm
{\bf 4.3}. ---

\vskip .3cm
{\bf 4.4}. ---

\vskip .3cm
{\bf 4.5}. ---


%%%%%%%%%%%%%%%%%%%%%%%%%%%%%%%%%%%%
\subsection*{5. Plongement cellulaires.}
\addcontentsline{toc}{subsection}{5. Plongement cellulaires}

{\bf 5.1}. --- 

\vskip .3cm
{\bf 5.2}. ---

\vskip .3cm
{\bf 5.3}. ---

\vskip .3cm
{\bf 5.4}. ---

\vskip .3cm
{\bf 5.5}. ---

%%%%%%%%%%%%%%%%%%%%%%%%%%%%%%%%%%%%
\subsection*{6. Plongements minimaux.}
\addcontentsline{toc}{subsection}{6. Plongements minimaux}

{\bf 6.1}. --- 

\vskip .3cm
{\bf 6.2}. ---

\vskip .3cm
{\bf 6.3}. ---


%%%%%%%%%%%%%%%%%%%%%%%%%%%%%%%%%%%%
\subsection*{7. Quelques remarques sur le cas des plongements dans les surfaces à bord non orientées (pouvant être non orientables).}
\addcontentsline{toc}{subsection}{7. Quelques remarques sur le cas des plongements dans les surfaces à bord non orientées (pouvant être non orientables)}

On peut introduire la notion de voisinage tubulaire standard non orienté (et éventuellement non orientable) d'un $1$-complexe $K$. Ces voisinages tubulaires sont classifiés pour la donnée d'un revêtement étale à deux feuillets $\tilde{K}$ de $K$ et d'une circularisation $\tilde{\omega}$ ``antisymétrique'' de $\tilde{K}$ (i.e.~: $\tau_* \tilde{\omega} = \tilde{\omega}^{-1}$, où $\tau$ est l'automorphisme involutif et sans point fixe du revêtement $\widetilde{K} \to K$). On peut aussi classifier ces voisinages tubulaires par des couples $(\omega, \xi)$ où $\omega$ est une circularisation de $K$, et $\xi$ une classe de cohomologie de $\mathrm{H}^1(K; \mathbf{Z}/2\mathbf{Z})$. On obtient alors un théorème de classification à congruence près (non orientée) des plongements dans les surfaces non orientées, théorème d'énoncé beaucoup plus technique que {\bf 3.4.3}. L'essentiel des résultats est donné dans [63].


% End
